\icode{E-MM-SI-L30-260625-003}
台球是一项深受青少年喜爱的体育运动。如图所示，水平台面上的甲球球心和中袋中心连线与桌边间夹角\(\theta =\) \,\qty{60}{\degree}，运动员球杆将白球乙沿与桌面平行的方向以速度\(v_0\)出击，乙与甲碰撞后甲落入中袋。已知两球质量、半径均相等，两球碰撞为弹性碰撞，不计两球所受阻力，不考虑球的自旋影响。关于甲、乙两球的运动，下列说法正确的是\paren[]

\begin{tikzpicture}
    \draw[dashed] (0,0)node[above,font=\tiny]{\kaishu 中袋}--(-60:{2*sqrt(3)/3})coordinate(a)--++(1,0)coordinate(b);
    \draw[fill=white] (-2,-.2) arc(-90:0:.2) (-1.8,-2) arc(0:90:.2) (1.8,-2) arc(-180:-270:.2) (2,-.2) arc(-90:-180:.2) (.2,-2)arc(0:180:.2) (.2,0)arc(0:-180:.2)  ($(a)+(120:.2)$)node[left,font=\tiny]{\kaishu 甲} circle(.1) (b) circle(.1)node[below,font=\tiny]{\kaishu 乙};
    \draw[fill=white,dashed] (a) circle(.1);
    \draw (-2,0) rectangle(2,-2);
    \draw[red,->] ($(b)+(-.1,0)$)--++(-.4,0)node[above,font=\scriptsize]{\(v_0\)};
\end{tikzpicture}
\begin{choices}
    \item 碰撞后乙球速度方向不变
    \item 碰撞后两球速度方向垂直
    \item 碰撞后甲球速度大小为\(\frac{\sqrt{3}}{2}v_0 \)
    \item 碰撞后乙球速度大小为\(\frac{v_0}{2}\)
\end{choices}
\begin{answer}
    B\quad 提示：
\end{answer}